\chapter{The Tendency}

\epigraph{\emph{There is something there that deserves a name.}}{}

This document begins from a deliberately modest observation. Across a set of disciplines that rarely speak to one another in a sustained way, researchers keep returning to a similar family of findings: structured relations among oscillatory, wave-bearing, and resonant processes seem to matter more than a purely arbitrary description of events would predict. Simple ratios, recurrent couplings, phase-locked regimes, interference patterns, and mode structures repeatedly appear as privileged constraints rather than as decorative byproducts. No complete theory yet coordinates these findings, and no single mechanism has been demonstrated everywhere. But a recognizable tendency has emerged across multiple literatures, and that tendency has not yet been named, coordinated, or investigated as a coherent field.

The breadth of that tendency is easiest to see when one begins not with music but with physics. Maxwell's field theory gave one of the decisive modern constructions of physical reality in terms of propagation, coupling, and wave-bearing fields rather than inert corpuscles alone \parencite{maxwell_1865}. De Broglie's wave mechanics and the electron diffraction experiments of Davisson and Germer extended that shift into the domain of matter itself, showing that the language of wavelength, interference, and periodic structure was not restricted to optics or acoustics \parencites{debroglie_1925}{davisson_1927}. More recent physics does not collapse into a single ontology, but it repeatedly preserves the same grammar: gravitational waves become observable events within a scientific construction of spacetime rather than speculative metaphor \parencite{abbott_2016}, and pattern formation in driven media appears across classical and quantum systems alike, from Faraday instabilities to Bose-Einstein condensates \parencites{engels_2007}{nguyen_2019}. Taken together, these developments make intelligible a research program that treats relation, mode structure, and resonance as candidates for explanatory priority.

The tendency is visible with particular clarity in neuroscience. Work on neural resonance has argued that tonal perception is not adequately described as passive registration of isolated frequencies, but rather as the dynamic entrainment of oscillatory systems to structured periodic relations \parencites{large_2012}{large_2005}. Research on cross-frequency coupling likewise suggests that simple phase-amplitude or phase-phase relations can operate as integrative mechanisms in the brain rather than as incidental byproducts of neural activity \parencite{canolty_2010}. More recent work on traveling waves and eigenfrequency structure extends this point: ratios among local frequencies may help coordinate geometric and computational organization across cortex \parencite{jacobs_2025}. These traditions do not converge on a single doctrine, but they repeatedly indicate that patterned relation is computationally and biologically salient.

Psychoacoustics and the study of vocal and musical structure constitute one especially legible branch of this larger tendency. Classical work on roughness and critical-band interaction framed consonance and dissonance in terms of interference patterns and spectral interaction rather than in terms of isolated tones understood atomistically \parencites{plomp_1965}{helmholtz_1863}. Later approaches complicated that picture by showing that auditory preference and tonal organization are also constrained by the statistics of vocal production and by the physiology of hearing \parencites{schwartz_2003}{bowling_2015}{bidelman_2009}. Parallel traditions built formal measures around periodicity, harmonic entropy, and spectral organization in order to explain why some intervallic relations are more stable, more compressible, or more perceptually coherent than others. The specific models differ, but the broader pattern is stable: harmonic relations are repeatedly treated as privileged constraints on perception and organization.

In ecology and biosemiotics, the same pressure reappears under yet another vocabulary. Krause's work on the acoustic niche hypothesis describes ecosystems as organized fields in which organisms partition spectral and temporal resources rather than emitting indiscriminately into noise \parencite{krause_2013}. The result is not simply coexistence, but patterned coexistence: recurrence, spacing, and non-random occupancy become ecologically meaningful. Biosemiotic traditions, meanwhile, have long argued that living systems do not merely exchange matter and energy; they exchange signs, and those signs depend on forms of recurrence and interpretability that can be stabilized across contexts \parencites{hoffmeyer_2008}{barbieri_2008}. What changes from one domain to another is the ontology of the object under study. What persists is the suspicion that patterned relation is not epiphenomenal. It does work.

The tendency is equally visible in nonlinear dynamics. Kuramoto's framework for coupled oscillators showed in paradigmatic form that simple rational relations are not arbitrary curiosities but privileged sites of phase locking and stable coordination \parencites{kuramoto_1984}{pikovsky_2001}. Arnold tongues, devil's staircases, and related synchronization structures make the point even more sharply: simpler ratios often command wider basins of locking and more robust coordination than more complex alternatives \parencite{strogatz_2000}. Similar mode-locking phenomena also appear in biological oscillators, where cardiorespiratory and other coupled rhythms can settle into stable integer relations rather than arbitrary drift \parencite{glass_1988}. Related work on recurrence has shown that what acoustics names as consonance can be formalized more generally as a measurable property of signal organization: simpler harmonic relations tend to generate stronger and more stable recurrence structure in the time domain \parencite{trulla_2018}. Even visualizations such as Lissajous curves make this point accessible in geometric form: ratio simplicity is expressed as topological simplicity and repeatable closure, not merely as aesthetic preference \parencite{gallozzi_2023}. Here again, the vocabulary differs from neuroscience or psychoacoustics, but the structural intuition is strikingly similar.

Comparative and cross-species work extends the question beyond the human case. Surveys of musicality across species, studies of consonance processing in non-human animals, work on pitch perception across species, and classic motivation-structural rules in vocal communication all suggest that sensitivity to intervallic, harmonic, or structurally patterned relations cannot be treated as an exclusively human cultural artifact \parencites{hoeschele_2015}{toro_2017}{walker_2019}{morton_1977}. The evidence is heterogeneous and should not be smoothed into a false consensus. But that heterogeneity is itself informative: once the question survives changes of species, niche, and task, it begins to look less like a local convention and more like a recurring problem in organized signaling.

At the outer edge of this survey, a smaller but increasingly explicit literature pushes the question further toward mind, matter, and consciousness. The Orch OR line developed by Hameroff and Penrose proposes that coherent quantum processes in microtubular structures may participate in conscious episodes, while Stapp's quantum interactive dualism reopens the causal status of conscious agency within quantum description \parencites{hameroff_2014}{stapp_2005}. Conscious realism and analytic idealism radicalize the displacement from another side by treating spacetime and matter as derivative appearances of a deeper field of conscious relation rather than as self-sufficient primitives \parencites{hoffman_2008}{kastrup_2017a}{kastrup_2017b}. HIT enters this frontier from a different side. It begins with cross-domain evidence of relational organization and only later asks how far the ontological consequences may reach. The existence of these peer-reviewed programs matters here because it shows that questions linking wave, field, matter, and consciousness are active sites of inquiry rather than mere cultural residue.

These literatures do not yet agree on a single explanatory center. What matters is that they keep rediscovering, from within their own methods and their own problems, that certain proportional relations and resonant organizations function as privileged sites of coupling, stability, compressibility, or interpretability. One field calls the phenomenon resonance. Another calls it periodicity. Another frames it as mode locking, wave propagation, harmonic entropy, interference structure, or recurrence. The names change because the local problems change. The structural drift does not.

Harmonic Information Theory therefore begins by naming a convergence problem. There appears to be a distributed empirical intuition, recurring across fields, that structured relations among oscillatory and wave-bearing processes carry or constrain information in ways that matter for physical stability, biological coordination, perception, and signal organization. Yet there is no Ariadne thread linking the literatures. Each field largely cites its own canonical authors, defines its own admissible evidence, and treats adjacent traditions as peripheral at best.

HIT is proposed first as an act of coordination. Its opening wager is disciplined: when multiple domains repeatedly discover that structured relations among oscillatory, resonant, and field-like processes matter, it becomes scientifically reasonable to ask whether those discoveries are shadows of a broader problem that has not yet been properly assembled. This wager does not reduce every important pattern in nature to simple ratios, nor flatten reality into a slogan about vibration. The acoustic and musical literatures matter here because they have developed one of the richest vocabularies for describing such relations. If a real pattern of convergence exists, the next problem is no longer merely empirical. It is epistemic.
